\documentclass[10pt,a4paper,ngerman,titlepage,tocindentauto]{scrartcl}

\usepackage{jhla4}
\usepackage[T1]{fontenc}
\usepackage[utf8x]{inputenc}
\usepackage{color}
\usepackage[colorlinks=true,breaklinks=true,linkcolor=darkblue,urlcolor=darkblue]{hyperref}
\usepackage{ngerman}
\usepackage{libertine}
\usepackage{graphicx}
\usepackage{tocstyle}
\usepackage{multicol}
\usepackage{listings}
\usepackage{wrapfig}
\usepackage{placeins}

\hypersetup{
	pdftitle={Software Projekt Uebersetzerbau 2010: Entwicklerdokumentation},
	pdfauthor={Rene Kijewski},
	pdfsubject={},
	pdfcreator={},
	pdfproducer={},
	pdfkeywords={},
	pdfkeywords={}
}

\definecolor{darkblue}{rgb}{0,0,0.5}
\definecolor{crimson}{rgb}{0.862745098,0.0784313725,0.235294118}

\usetocstyle{allwithdot}
\setcounter{tocdepth}{3}
\setcounter{secnumdepth}{3}

\renewcommand{\listfigurename}{Abbildungsverzeichnis}
\renewcommand{\lstlistingname}{Quelltext}
\renewcommand{\lstlistlistingname}{Quelltextverzeichnis}

\newcommand{\xmlNote}[2]{\footnote{Spezifikation zu \href{http://www.w3.org/TR/2008/REC-xml-20081126/\#{#1}}{\texttt{#2}}.}}
\newcommand{\link}[1]{\href{#1}{#1}}
\newcommand{\TODO}{ {\LARGE\bf\color{crimson} TODO} }
\newcommand{\n}{\linebreak[1]}
\newcommand{\java}[1]{\lstinline[language=java,showstringspaces=true]{#1}}

\begin{document}
	\begin{titlepage}
		\titlehead{\includegraphics[width=\textwidth]{graphiken/Fub-logo.pdf}}
		\title{Softwareprojekt Übersetzerbau 2010} 
		\subtitle{Entwicklerdokumentation}
		\author{René Kijewski} 
		\date{18. Juli 2010 \vfill}
		\maketitle
	\end{titlepage}
	
	\begin{abstract}
		Das Softwareprojekt Übersetzerbau 2010 setzte sich zum Ziel, einen XML- nach Java"=Source"=Code"=Übersetzter
		zu schreiben.
		Dieses Dokument erklärt im Speziellen die Komponenten XML"=Lexer, Statement"=Lexer
		und die Symboltabelle. Es richtet sich an Entwickler und enthält deshalb gegebenenfalls
		Jargon und versucht auch nicht, Sachverhalte zu vereinfachen.
		Die entsprechende Kundenpräsentation, die sich nicht an Programmierer richtet, steht unter
		\link{https://dev.spline.de/svn/ss10-swp-uebersetzerbau/trunk/dokumentation/rk-kundenpraesentation/rk-kundenpraesentation.pdf}
		zur Verfügung.
	\end{abstract}

	{
		\pagestyle{empty}
		\tableofcontents
		\newpage
	}
	
	\parskip 7pt plus 5pt minus 3pt
	\setcounter{page}{1}
	\section{Einleitung}
		Das Softwareprojekt Übersetzerbau 2010 machte es sich zur Aufgabe, eine selbst"=erdachte Programmiersprache
		zu spezifizieren und einen Compiler für diese zu Implementieren.
		
		Die Projektgruppe, bestehend aus -- in alphabetischer Ordnung --
			Alexander Rau,
			Ansgar Schneider,
			Igor Merkulow,
			Markus Rudolph,
			René Kijewski und
			Stefan Meißner,
		entschied sich, eine XML"=artige Eingabesprache zu benutzen, die
		\href{http://java.sun.com/javase/6/}{Java"=SE"=6}"=kompatiblen Sourcecode ausgibt.
		Dementsprechend wurde auch Java 1.6 als Programmiersprache gewählt.
		Als einheitliche Programmierungebung wurde \href{http://www.eclipse.org/galileo/}{Eclipse Galileo} verwandt.
		Apache \href{http://subversion.apache.org/}{Subversion}, gehostet bei \href{http://dev.spline.de/}{Spline}, wurde als
		Versionskontrollsystem benutzt. Der Quelltext kann von \link{https://dev.spline.de/svn/ss10-swp-uebersetzerbau/}
		heruntergeladen werden.
		
		Ich übernahm die Komponenten XML"= und Statement"=Lexer, sowie die Implementierung einer Symboltabelle.
		
		Bei den von mir implementierten Teilprojekten achtete ich auf hohe Wiederverwendbarkeit
		und versuchte, eine Benutzung besonders einfach zu gestalten. So legte ich größeren Wert darauf,
		dass Weiternutzer die Komponenten ohne lange Einlesezeit benutzen können als auf besonders große
		Performanz. Dieser Ansatz musste leider später aufgewicht werden. In den Abschnitten
		\hyperlink{Arbeitsverlauf}{Arbeitsverlauf} und \hyperlink{Abschlussbemerkungen}{Abschlussbemerkungen}
		finden Sie weitere Informationen dazu.
		
		Um den Einstieg in die Nutzung zu Vereinfachen, verwandte ich ausgiebig die verbreitete Entwurfsmuster,
		vorrangig das \href{http://c2.com/cgi/wiki?IteratorPattern}{\em Iterator"=Patter} und die
		\href{http://c2.com/cgi/wiki?FactoryMethodPattern}{\em Factory"=Methode}.
		Beide Pattern werden im Buch {\em Design Patterns: Elements of Reusable Object"=Oriented Software}%
		\footnote{ISBN \href{https://portal.d-nb.de/opac.htm?method=simpleSearch&query=0201633612}{0-201-63361-2}}
		beschrieben und sind vielen Entwicklern von objektorientiertem Code geläufig.
		
		Die einzelnen Komponenten trennen so weit wie möglich zwischen Interface und Implementierung.
		Dies sollte es bei einer späteren Weiterarbeit am Projekt vereinfachen, die Komponenten
		auszutauschen. Bei der Spezifikationsphase überlegten wir uns unter anderem, dass auch eine
		Compilierung in Java"=Bytecode nativ möglich sein sollte.%
		\footnote{Das heißt, dass wir nicht von \texttt{javac} abhängig sein wollten.}
		Zudem sollte auch die Ausgabesprache auswechselbar sein, so dass nur das Kernstück,
		der Syntaxbaum, sowie das Interface der weiteren Komponenten, also Symboltabelle und Builder,
		unveränderlich wäre.
		
	\section{Teilprojekte}
		Die einzelnen Teilprojekte sind für sich allein verwendbar -- die Verzahnung
		mit anderen Komponenten ist minimal --, da sie jeweils am Anfang der Bearbeitungskette stehen.
		
		\begin{figure}[ht]
			\caption[Verlaufdiagramm]
				{\hypertarget{Verlaufdiagramm}{Verlaufdiagramm (in diesem Dokument beschriebene Komponenten hervorgehoben)}}
			\fbox{\centering
				\includegraphics[width=\textwidth]{graphiken/flowchart.pdf}
			}
		\end{figure}
		
		Die Lexer"=Komponenten benutzen jeweils einen Zeichenstrom als Eingabe. Die Rückgabedaten speichern die
		Stelle im Eingabestrom an der sie standen, so dass spätere Komponenten, wenn sie einen Fehler
		in der Eingabe feststellen, eine geeignete Meldung samt Position ausgeben können.
		Die Symboltabelle verwendet dasselbe Interface, um die Eingabeposition zu speichern, jedoch ist
		die Bedeutung an dieser Stelle wahrscheinlich geringer als bei den Lexern.
		
		\subsection[XML"=Lexer]{\hypertarget{Teilprojekte_XML_Lexer}{XML"=Lexer}}
			Der XML"=Lexer {\em xmlNodeStream} vereinfacht die Arbeit des DOM"=Parsers, indem er die syntaxtischen
			Komponenten des Eingabestromes auftrennt, das heißt {\em tokenized}.
			Die erkannten Komponenten sind Textknoten, Kommentare, Tags, Endtags, Attribute und Processing instructions.
			
			Als Eingabeformat wird XML 1.0 verwendet.%
			\footnote{\href{http://www.w3.org/TR/2008/REC-xml-20081126/}{Extensible Markup Language (XML) 1.0 (Fifth Edition)}}
			Es wurden bewusst Eineinfachungen jedenüber dem W3-Standard vorgenommen. Unter anderem sind
			nicht wohlgeformte Ausdrücke nach Art ,,\verb|<tag attr="a < b"/>|`` erlaubt.
			Dieses Beispiel wäre aufgrund des Kleiner"=Als"=Zeichens im Attributewert
			nach der Spezifikation verboten.\xmlNote{NT-AttValue}{AttValue}
			
			Die API des {\em xmlNodeStreams} ist an {\em StAX}, der \href{http://stax.codehaus.org/}{Streaming API for XML},
			angelehnt. StAX hat sich seit der Spezifikation im \href{http://jcp.org/en/jsr/detail?id=173}{JSR 173} aus
			dem Jahr 2006 als Standard für einfaches Parsen von XML"=Eingaben etabliert.
			xmlNodeStream besitzt nicht dieselbe Mächtigkeit wie StAX und erfüllt in Folge dessen
			nicht die komplette Spezifikation des W3-Consortiums, da der XML"=Prolog\xmlNote{xmldoc}{prolog}
			nicht als solcher verstanden wird:
			\begin{itemize}
				\item
					Die XML"=Deklaration\xmlNote{NT-XMLDecl}{XMLDecl} wird als
					Processing instruction\xmlNote{NT-PI}{PI} verstanden.
				\item
					Die Doctype"=Deklaration\xmlNote{NT-doctypedecl}{doctypedecl} wird nicht unterstützt.
			\end{itemize}
			Konkretere Informationen über den Informationsfluss stehen in der \hyperlink{XML-Lexer-API}{API"=Referenz}.
			
			Der XML"=Lexer überprüft die Eingabe nicht auf Wohlgeformtheit, das heißt, auch
			,,\verb|<a></b>|`` würde nicht nicht als Fehler erkannt. Dies ist eine bewusste Designentscheidung,
			die einerseits den Overhead durch eine solche Überprüfung verhindern soll -- der DOM"=Parser
			führt eh über eine Überprüfung durch -- und andererseits die Anwendungsmöglichkeiten der Implementierung
			verbreitern soll: Gewöhnliche Webseiten sind selten wohlgeformtes XML. Dadurch,
			dass {\em xmlNodeStream} agnostisch gegenüber Strukturfehlern ist, wäre er auch zum
			Lexen von Internetseiten geeignet. Fehler wie fehlende abschließende Anführungszeichen werden
			jedoch nicht übergangen, da sie eine aktive Fehlerkorrektur erforderten.
			
			{\em xmlNodeStream} implementiert das Iterator"=Interface Javas, wodurch der Einstieg in die
			Benutzung vereinfacht werden soll.
			
			\subsubsection*{Beispiel}
				\begin{lstlisting}[frame=single,language=XML,caption=Beispiel.xml]
<?xml version="1.0"?>
<module name="Test">
	Inhalt
</module>
				\end{lstlisting}
				Die Eingabedatei obenstehende XML"=Datei würde folgendermaßen erkannt:
				\begin{itemize}
					\item Processing instruction mit dem Namen \verb|xml| und Wert \verb|version="1.0"|.
					\item Öffendes Tag mit dem Namen \verb|module|.
					\item Attribute mit dem Namen \verb|name| mit dem Wert \verb|Test|.
					\item Textknoten mit dem Wert \verb|Inhalt|.
					\item Schließendes Tag mit dem Namen \verb|module|.
				\end{itemize}
				An diesem Beispiel Erkennt man, dass das jedes eingelesene Token ein Tripel aus
				Typ, Name und Wert ist, wobei Name und/oder Wert auch leer beziehungsweise \texttt{null} sein können.
			
			\begin{figure}[ht]
				\caption[Übergangstabelle der Eingabedaten]{\hypertarget{Uebergangsdiagramm_XML_Lexer}{Übergangstabelle der Eingabedaten}}
				\fbox{\centering
					\includegraphics[width=1.4\textwidth,angle=90]{graphiken/xml.pdf}
				}
			\end{figure}
			
		\subsection{Statement"=Lexer}
			Der Statement"=Lexer hat eine ähnliche Aufgabe wie der XML"=Lexer: Einen Eingabestrom in seine
			Komponenten auftrennen. Das Lexen der Statements ist die Vorstufe zum Parsen.
			
			Der Statement"=Lexer verwendet das Iterator"=Interface, wobei die Rückgabedaten
			Tupel aus dem Typen des Datums und eventuell einem zugehörigen Wert, zum Beispiel bei Strings.
			In der derzeitigen Implementierung gibt es 33 Arten von erkannten Token, sowie zwei
			Statuscodes für das Eingabeende beziehungsweise bei lexikalisch fehlerhaften Eingaben.
			Der Umfang der erkannten Eingabefehler ist gering; die Fehlererkennung geschieht im
			{\em Statement"=Parser}. Als Fehler würde unter anderem ,,\verb|1a|`` erkannt; ,,\verb|a b c|`` hingegen nicht.
			
			Die Syntax der Statements lehnt sich an Java\,SE\,6 an. Da Javas Operatoren nur ein oder
			zwei Zeichen lang sind\footnote{Der Operator \texttt{{>}>{>}} wird nicht verstanden.} und es keine Zustände wie beim
			Lexen von XML"=Daten wird, wird keine Übergangstabelle benutzt, sondern ein \verb|switch|"= und
			\verb|if|"=Konstrukt. Dies vereinfacht im Allgemeinen sowohl das Lesen und Schreiben des Lexers,
			erschwert jedoch Spezialfälle. So wird zum Beispiel zum Erkennen der Token \verb|new| und
			\verb|null| auf Tricks wie Push"=Back"=Buffer zurückgegriffen.
			Ein entscheidender Vorteil ist jedoch, dass das Einlesen der Daten wesentlich schneller
			als beim XML"=Lexer geschieht.
	
		\subsection{Symboltabelle}
			\TODO
	
	\section[Arbeitsverlauf]{\hypertarget{Arbeitsverlauf}{Arbeitsverlauf}}
		Die Programmierung des XML"=Lexers wurde bereits am 30. April begonnen; die Komponente ist deshalb
		wohl am ausführlichsten getestet. Die Aufgabe des Statement"=Lexers wiederum wurde
		am 19. Mai begonnen und die Arbeit an der Symboltabelle wurde am 28. Mai aufgenommen,
		wudurch gerade für die letzte Aufgabe ein erheblicher Zeitdruck bestand.
		
		\subsection{XML"=Lexer}
			Die Arbeit am XML"=Lexer verlief verhältnismäßig zielgerichtet. In der Vorlesung
			Übersetzerbau wurde ausgiebig erklärt, wie ein Parser geschrieben wird, so dass
			sich die tatsächliche Umsetzung als einfach erwies. Probleme ergaben sich nur,
			wenn ich Übergänge im Zustandsdiagramm (siehe \hyperlink{Uebergangsdiagramm_XML_Lexer}{Abbildung 2})
			übersah, weshalb ich sehr schnell dazu überging, das Diagramm in der Sprache
			\href{http://www.graphviz.org/doc/info/lang.html}{DOT} festzuhalten. Zeichnungen auf
			dem Papier erwiesen sich nicht als sachdienlich.
		
			\TODO
		
		\subsection{Statement"=Lexer}
			Der Statement"=Lexer wurde in kurzer Zeit geschrieben.
			
			\TODO
		
		\subsection{Symboltabelle}
			Die Arbeit an der Symboltabelle erwies sind als kompliziert. Da die Aufgabe leider sehr spät erst
			vergeben wurde und aufgrund der Anforderung, dass sie Typisierung beherrschen sollte, gehörte
			die Symboltabelle planerisch zu den größten Komponenten.
			
			Die Symboltabelle sollte zunächst für sämtlichen Zielsprachen gleich implementiert sein. Nachdem
			ich mich jedoch weiter in die Thematik einlas, erwies sich das nicht als möglich, da die verschiedenen
			Zielsysteme zu unterschiedlich sind. So haben verschiedene Programmiersprachen, selbst wenn man
			sich auf objektorientierte Programmiersprachen beschränkt, wenig Gemeinsamkeiten. Als Beispiele gegenüber
			Java sein die funktionale Programmiersprache \href{http://caml.inria.fr/ocaml/}{OCaml} und die prototypen"=basiere
			Programmiersprache \href{http://www.ecma-international.org/publications/standards/Ecma-262.htm}{ECMAScript}
			angeführt, deren Klassenkonzept keine Gemeinsamkeiten mit Java aufweisen, jedoch sinnvolle Zielsprachen
			wären.
			
			Ich beschränkte mich auf die Umsetzung einer Symboltabelle, die Java 1.6 sowohl für die Ausgabe
			nach Quelltext als auch Bytecode beherrt. Das Interface wiederum sollte so generisch gehalten werden,
			dass später weitere Zielsysteme eingeführt werden könnten, wobei sich ein zu generischen Interface
			als zu kompliziert -- und damit als unbenutzbar -- erwies. Herr Rudolf gab mir des Öfteren hinweise
			auf Verbesserungsmöglichkeiten am Interface, die mir zuvor nicht einfielen. Wiederum war das Interface
			dadurch während der gesamten Entwicklungszeit einem ständigen Wandel unterlegen, wobei ich mir eigentlich
			eine schnelle Festlegung wünschte, gegen die ich die konkrete Implementierung entwickeln könnte.
			Retrospektiv kann ich sagen, dass ich mich anfangs zu wenig mit meinen anderen Teammitgliedern
			zusammengesetzt habe und deshalb nicht die Fülle der Anforderungen kannte.
			
			Bei der Programmierung an der Symboltabelle habe ich einige wertvolle Kenntnisse über die Interna
			Javas gewonnen. Gerade bei der Reflexion"=API, die einem Programm zur Laufzeit Informationen über
			das Klassen"= beziehungsweise Typsystem geben soll, habe ich erhebliche Wissensprünge gemacht,
			die ich auch für andere Kurse nutzen konnte.\footnote{Siehe zum Beispiel \href
			{https://dev.spline.de/svn/xmlproject10-group-x/trunk/controller/src/java/de/spline/dev/xmlproject10_group_x/imports/xmlContainers/Container.java}
			{\texttt{Containers.java}} des XML"=Technologien"=Projekts} Leider erwies sich die sehr einfach zu
			benutzende Reflexions"=API als Haupt"=Flaschenhals des gesamten Projektes, weshalb ich einige ,,halbherzige``
			Behelfsmaßnahmen eingebaut habe -- {\em halbherzig} deshalb, weil ich das Problem nur versteckt und
			nicht beseitigt habe. Die Reflexions"=API ist leider äußerst langsam und verbraucht große Mengen an
			RAM"=Speicher, da die Inspizierten Klassen, sowie sämtliche (transitiven) Abhängigkeiten in den Speicher
			geladen werden, wobei eigentlich nur die Struktur der Klasse bedeutsam für die Erfüllung der Aufgabe wäre.
			
			Als erste Behelfsmaßnahme werden nur noch Klassen in die Symboltabelle gelesen, die auch von außen sichtbar
			sind, das heißt entweder \texttt{public} oder \texttt{protected} sind. Default"=Private und \texttt{private}"=Klassen
			werden nicht mehr eingelesen. Jedoch offenbart diese Beschreibung bereits das größte Problem dieser Umsetzung:
			Nur weil Klassen nicht adressierbar {\em wäre}, gilt das nicht für zum Beispiel Rückgabewerte. Als weitere
			Maßnahme zum Beschleunigen habe ich das Laden pa\-ral\-le\-li\-siert, wodurch sich -- gefühlt -- ein fast linearer Anstieg
			der Geschwindigkeit mit der Zahl der Benutzen Rechenkernen ergab.
	
	\section{Offene Baustellen}
		Leider blieben aufgrund der geringen Projektdauer noch einige Komponenten fehlerhaft oder unvollständig.
		Jene Verbesserungsvorschläge, die mir einfielen, hielt ich im \href
		{https://dev.spline.de/svn/ss10-swp-uebersetzerbau/trunk/dokumentation/Offene\%20Baustellen.txt}
		{\texttt{Offene Baustellen.txt}} beziehungsweise als \texttt{TODO}s an der Entsprechenden Stelle im
		Quelltext fest. Die einzelnen Teilprojekte sind noch nicht ausreichend durch Unittests überprüft worden.
		Bei einer Weiterarbeit an dem Projekt täten folgende Programmierer gut daran, zuerst Blackbox-Text
		der einzelnen Komponenten durchzuführen, um eventuelle Verstöße gegen das Interface festzustellen.
	
		\subsection{XML"=Lexer}
			Der XML"=Lexer steht an erster Stelle in der des \hyperlink{Verlaufdiagramm}{Arbeitsablaufes} und kann deshalb,
			trotz des Fehlens formeller Unittests als überprüft betrachtet werden. Eventuelle Fehler hätten sich
			beim Gebrauch zeigen sollen.
			
			Offen ist noch die komplette Implementierung des XML"=1.0"=Standards, wie bereits
			\hyperlink{Teilprojekte_XML_Lexer}{oben erwähnt}, wobei meines Erachtens der Mehrwert
			durch die vollständige Implementierung des Standars nicht der Aufwand rechtfertigen
			würde.
	
		\subsection{Statement"=Lexer}
			Der Statement"=Lexer enthält noch einen dokumentierten Fehler, dessen Ursache noch nicht ausgemacht werden konnte.
			Der Bezeichner ,,\texttt{nul}`` ist nicht erlaubt, was mit der Implementierung der Erkennung des Keywords
			\texttt{null} zusammenhängen wird. Das Erkennen der Schlüsselwörter ist derzeit suboptimal implementiert, da
			mir leider bei bei der Arbeit an der Symboltabelle ein besserer Algorithmus einfiel.
			
			Derzeit werden wird die Eingabe mit einem Pushback"=Buffer eingelesen. Sollten Zeichen gelesen werden,
			die nicht zu dem derzeitigen Token gehören, werden sie wieder in den Eingabestrom zurückgeschrieben.
			Sollte nur ein Zeichen zurückgeschrieben werden müssen, arbeitet diese Methode performant und ist zudem
			einfach zu implementieren. Sollten jedoch mehr Zeichen zurückzuschreiben sein, so kann sich diese Methode
			als kompliziert zu warten erweisen. Zuerst schrieb ich einen Pushback"=Buffer mit einem Zeichen als
			Pushback und schrieb ihn im weiteren Arbeitsverlauf zu einem allgemeinen Puffer um, in dem man beliebig
			viele Zeichen zurückschreiben kann. Vermutlich ist mit dabei ein Fehler unterlaufen, der in den Beobachtung
			resultiert, dass der Bezeichner ,,\texttt{nul}`` nicht erlaubt ist.
			
			Eine sinnvollere Implementierung als die gleichzeitige Überprüfung, ob in der Eingabe ein Identifier oder
			ein Schlüsselwort steht, wäre, dass zuerst nur Identifier gelesen werden und, nachdem ein vollständiges
			Token gelesen wurde, überprüft wird, ob dieses zu einem Schlüsselwort wie \texttt{null} gehört.
	
		\subsection[Symboltabelle]{\hypertarget{Offene_Baustellen_Symboltabelle}{Symboltabelle}}
			Im weiteren Programmier"= und Testverlauf erwies sich, dass die Symboltabelle optimiert werden muss, um
			nutzbar zu sein, da die Ladezeit auf Windows"=Laptops mehrere Minuten betrug. So machte es zum Beispiel -- und zwar in einem Maße,
			das mich selbst erstaunt hat -- einen Unterschied, ob eine Stringeingabe gegen $10$ oder $\log 10$ Strings verglichen
			wird, letzteres durch die Verwendung eines \verb|Set<String>|. Häufig ist der iterative Vergleich von von kleinen
			Mengen von Werten, hier zirka zehn, schneller als ein optimierter, der in $\textsf O\!\left(\log n\right)$ arbeitet,
			jedoch scheint der Vergleich von Strings in Java um mehrere Zehnerpotenzen langsamer als bei C zu sein.
			
			Um das Problem der hohen Ladezeit zu umgehen und den Speicherbedarf minimal zu halten wären zwei Ansätze,
			idealer Weise kombiniert, gangbar.
			\begin{itemize}
				\item
					Die Reflexions"=API ist zu mächtig und schwerfällig. Ein manuelles Auslesen der klassenstruktur
					auf Objekt"=Code"=Ebene ließe es zu, dass man unnötige Informationen nicht im Speicher hält beziehungsweise
					gar nicht erst einliest.
				\item
					Ein Nadelöhr stellt die
					Benutzung von \texttt{TreeMap}s dar, die ich an vielen Stellen verwende, da sie im Gegensatz zu \texttt{HashMap}s
					nicht die Gefahr von Kollisionen in sich bergen. Vermutlich ist die Datenmenge klein genug, so dass es zu keinen
					Kollisionen der Streuwerte käme. Eine Überprüfung dieser Annahme wäre jedoch erforderlich.
			\end{itemize}
	
	\section[Abschlussbemerkungen]{\hypertarget{Abschlussbemerkungen}{Abschlussbemerkungen}}
		Eine Lexion des Projektes war, dass Besprechungen notwendig sind. Noch bevor man eine Idee umsetzt, sollte man
		sich mit den betroffenen Teammitgliedern zusammen setzten und um gemeinsam die Implikationen abzuschätzen
		oder Verbesserungsvorschläge einholen zu können. Ebenso lassen sich Beschreibungen wesentlich besser vis"=à"=vis
		vermitteln als das mit Texten möglich wäre. Als Programmierer ist man häufig zu tief in der Materie drin, als dass
		man sinnvolle Fragen abschätzen könnte, die in eine Dokumentation zu schreiben wären, wodurch zudem viel Zeit
		mit dem Schreiben unnötiger Kommentare verloren geht. Ich denke, dass Pair"=Programming das Mittel der Wahl ist,
		wenn es sich zeitlich organisieren lässt. Die Kommunikation über Chat oder Telefon kann leider nicht dasselbe leisten.
		
		Die wöchentlichen Teambesprechungen habe ich als reine Zeitverschwendung empfunden, da wöchentliche
		Berichte gerade in der Anfangsphase eines Projektes -- über die das Projekt aufgrund des Zeitdrucks
		auch nicht heraus kam -- viel zu selten sind und man zu viele für einen selbst unrelevante Informationen
		anhören muss, so dass das Aufmerksamkeitsfenster bei relevanten Dingen bereits wieder geschlossen
		sein kann. Dass die Teambesprechungen in unserem Fall um acht Uhr morgens stattfanden, war der Aufmerksamkeit
		ebenso nicht zuträglich.
		
		Nichtsdestoweniger hat das Projekt Spaß gemacht, wobei leider einige -- meines Erachtens eklatante -- Fehler in
		der Planung gemacht wurden. So haben wir viel zu viel Zeit damit aufgewendet, Lexer und Parser zu schreiben, wobei
		es zu diesem Zwecke bereits viele gute Open"=Source Tools gäbe, die wir hätten stattdessen verwenden sollen.
		Tatsächlich habe ich für die Verwendung einer XML"=ähnlichen Sprache optiert, damit wir einen vorhandene
		DOM"=Parser hätten benutzen können.
		
		Ich für meinen Teil hatte mir erhofft, mehr über die Optimierung von Zwischencode zu lernen; andere Projektmitglieder
		wiederum wollten eine Programmiersprache mit neuartigen Konzepten schaffen, wobei ich gerade Herrn Rudolfs Konzept einer
		Zustandsmaschine interessant fand. Ebenso hörte sich das Konzept einer einsteigerfreundlichen Programmiersprache gut an.\footnote
		{Ich muss gestehen, dass ich den Namen des Diplomstudenten vergessen habe.}
		Retrospektiv muss ich sagen, dass wir uns als Gruppe zu wenig Zeit genommmen haben, uns über die Möglichkeiten
		klarzuwerden. Leider haben wir die Konzeptphase zu schnell abschließen wollen, wobei ich jedoch hoffe, so einen
		Fehler in Zukunft nicht mehr zu machen.
		
		Zur Eskalation am Ende des Projektes hat meinerseits gerade geführt, dass das Projekt nicht die Dinge umfasste,
		die wir uns am Anfang vorgestellt haben. Ich wollte Probleme offen ansprechen, um sie eventuell besser angehen zu
		können -- Probleme, die nicht angesprochen werden, werden sich nur vergrößern. Leider ist mein Ton wohl zu aggressiv
		rübergekommen, wodurch ich jene Problem wohl nur noch verstärkt habe.
		
		Ferner hat jedoch die strikte Trennung zwischen ,,Management`` und den Programmieren zur Eskalation beigetragen
		oder hat sie erst möglich gemacht. Es bestand ein großes Misstrauen zwischen mir und dem ,,Management``,
		unter anderem geschuldet dadurch, dass uns nicht mitgeteilt wurde, was in den Wochenbesprechungen zwischen den
		Teamleitern und Frau Dr. Fehr beredet wurde, sowie dass uns Programmierern die Dokumentationen und Berichte
		nicht ausgehängigt wurde, als wir danach fragten.
		
		Ich denke, dass es für weitere Softwareprojekte wichtig ist, die Teamleitung dazu anzuhalten, sich auch in die
		Programmierung einzubringen, da ansonsten die Stimmung {\em unweigerlich} kippen {\em muss}. Dennoch denke ich,
		dass ich durch dieses Softwareprojekt eine Erfahrung fürs Leben gemacht habe --- mehr als es ein Besuch der
		\href{https://www.inf.fu-berlin.de/w/SE/VorlesungSoftwaretechnik2008}{Vorlesung Softwaretechnik} je könnte.
		
	\section{API"=Referenz}
		\subsection[XML-Lexer]{\hypertarget{XML-Lexer-API}{XML"=Lexer}}
			Der XML"=Lexer trennt zwischen Interface und Implementierung.
			Das Package des Interfaces heißt \texttt{de.{\n}fu\_berlin.{\n}compilerbau.{\n}xmlNodeStream};
			die Implementierung befindet sich unter \texttt{de.{\n}fu\_berlin.{\n}compilerbau.{\n}xmlNodeStream.{\n}impl}.
			
			\subsubsection*{Interface}
				\begin{wrapfigure}{R}{0.3\textwidth}
					\caption{Mög\-li\-che Rück\-ga\-be\-wer\-te des State\-ment"=Le\-xers nach Zu\-stand}
					\fbox{\centering
						\includegraphics[width=0.3\textwidth]{graphiken/Ausgaben.pdf}
					}
				\end{wrapfigure}
				
				Das Interface defininiert in \texttt{NodeType.java} die Arten der zu erkennenden Knoten.
				Dies sind derzeit \texttt{NT\_TEXT}, \texttt{NT\_COMMENT}, \texttt{NT\_TAG}, \texttt{NT\_END\_TAG},
				\texttt{NT\_ATTR}, \texttt{NT\_PI}, sowie für fehlerhaften Eingaben \texttt{NT\_ERROR}.
				
				Ein einzelnes gelesenes Token wird in einem \texttt{XmlNode} zurückgegeben, welches die Methoden
				\texttt{NodeType getType()}, \texttt{PositionString getKey()} und \texttt{PositionString getValue()}
				besitzt. Nicht für alle Tokenarten sind sowohl Key also auch Value gesetzt. Konsultieren dafür
				bitte das entsprechende JavaDoc für \texttt{XmlNode}. Ein \texttt{XmlNode} implementiert das Interface
				\texttt{StreamPosition}.
				
				Die letzte Komponente des Interfaces des XML"=Lexers stellt \texttt{XmlNodeStream} dar, welches
				angibt, dass die XmlNodeStreams das Interface \texttt{Iterable<{\n}XmlNode>}, \texttt{StreamPosition} und
				\texttt{Closeable} erfüllen müssen.
				
				Wie in Abbildung 3 zu sehen, kann der XmlNodeStream nicht an jeder Stelle ein beliebiges Token zurückgeben.
				So kann unter anderem \texttt{NT\_ATTR} nur auf \texttt{NT\_ATTR} oder \texttt{NT\_TAG} folgen.
				Bei der Benutzung von \texttt{XmlNodeStreamFactory} darf man sich darauf verlassen, dass dieser Contract
				nicht verletzt wird.
			
			\subsubsection*{Implementierung}
				Die Implementierung verfügt nur über eine öffentlich sichtbare Klasse, nämlich \texttt{XmlNodeStreamFactory}.
				Die Klasse verfügt über drei statische Factory"=Methoden
				\begin{itemize}
					\item \texttt{static XmlNodeStream createNewInstance(PositionString str);}
					\item \texttt{static XmlNodeStream createNewInstance(Reader reader);}
					\item \texttt{static XmlNodeStream createNewInstance(Reader reader, StreamPosition pos);}
				\end{itemize}
				Wobei die ersten beiden Methoden Wrapper für die letzte Methode darstellen.
	
		\subsection{Statement"=Lexer}
			Der Statement"=Lexer trennt zwischen Interface und Implementierung.
			Das Package des Interfaces heißt \texttt{de.{\n}fu\_berlin.{\n}compilerbau.{\n}statementLexer};
			die Implementierung befindet sich unter \texttt{de.{\n}fu\_berlin.{\n}compilerbau.{\n}statementLexer.{\n}impl}.
			
			Das Interface definiert in \texttt{TokenType.java} fünfunddreißig verschiedene Arten von gelesenen Tokens,
			wobei \texttt{EOF} und \texttt{ERROR} Statusmeldungen darstellen. Die Tokenarten orientieren sich
			an der Grammatik Javas und enthalten auch keine eigenen Elemente. Jedes gelesene Token wird in einem
			\texttt{StatementNode} zurückgegeben, wobei ein Node ein Tupel aus dem Typen und optional einem Wert
			darstellt.
			Tokens tragen nur dann einen Wert, wenn sie vom Typ \texttt{ID}, \texttt{STRING}, \texttt{INT} oder
			\texttt{REAL} sind. Zum Auslesen der Werte gibt es die Methoden
			\begin{itemize}
				\item \texttt{Object getValue() throws IllegalAccessException;}
				\item \texttt{Number getNumber() throws IllegalAccessException;}
				\item \texttt{CharSequence getString() throws IllegalAccessException;}
			\end{itemize}
			Eine \texttt{IllegalAccessException} wird geworfen, wenn das Token keinen (entsprechenden)
			Wert enthält. \texttt{getNumber()} enthält nur für \texttt{INT} und \texttt{REAL} einen Wert;
			\texttt{getString()} wiederum nur für \texttt{ID} und \texttt{STRING}.
			
			Das Tokenisieren geschieht über die Factory"=Methode \texttt{StatementLexer.tokenize(PositionCharacterStream)},
			wobei der Rückgabewert ein \texttt{Iterable<StatementNode>} ist und somit über die Methoden \texttt{next()} und
			\texttt{hasNext()} angesprochen werden kann.
	
		\subsection{Symboltabelle}
			Das Interface der Symboltabelle liegt im Package \texttt{de.{\n}fu\_berlin.{\n}compilerbau.{\n}symbolTable}.
			Die Implementierung ist -- wie der eng verbundene Annotator -- abhängig vom Zielsystem.\footnote
			{Angemerkt sei hierbei, dass die Implementierungen des Annotators und der Symboltabelle nicht von
			der konkreten anderen Implementierung abhängig sind.}
			Meine Implementierung der Symboltabelle hat als Zielsystem Java 1.6. Die konkrete Implementierung
			liegt in \texttt{de.{\n}fu\_berlin.{\n}compilerbau.{\n}symbolTable.{\n}java}. Die Ausnahmebehandlungsklassen
			liegen in \texttt{de.{\n}fu\_berlin.{\n}compilerbau.{\n}symbolTable.{\n}exceptions}. Die Exceptions sind
			ihrerseits unabhängig vom Zielsystem.
			
			Die Aufbau der Symboltabelle ist verhältnismäßig komplex. Um die Benutzung zu vereinfachen, habe ich
			auf ein Factory"=Pattern gesetzt. Aus dem Paket der Implementierung ist nur die Klasse von \texttt{RuntimeFactory}
			von außen sichtbar. Sämtliche anderen Klassen sind package-private. Diese Klasse wiederum hat nur eine
			public Methode:
			\begin{itemize}
				\item \verb|static Runtime newRuntime(URL[] classpath, URL rtJar)|
			\end{itemize}
			Mit dem Parameter \texttt{classpath} kann eine Liste mit JARs übergeben werden, deren Interface bei der Kompilierung
			verfügbar sein soll. Die Liste entspricht dem Parameter \texttt{-classpath} von \texttt{javac}.
			Die \texttt{rt.jar} wiederum enthält das Basissystem Javas, vornehmlich die Klassen in \texttt{java.lang}.
			Aufgrund der Benutzung der Reflexions"=API (siehe Abschnitt unter ,,\hyperlink{Offene_Baustellen_Symboltabelle}{Offene Baustellen}``)
			muss die übergebene \texttt{rt.jar} dieselbe sein, die die JVM des Compilers verwendet.
			Dieses Problem durch die Abkehr von der Reflexions"=API behoben werden.
			
			\begin{wrapfigure}{R}{0.5\textwidth}
				\caption[Vererbungshierarchie der Symbole]{Vererbungshierarchie der Symbole: \newline
					Factories gelblich hinterlegt \newline
					Metaklassen bläulich hinterlegt \newline
					qualified Symbols grün umrandet}
				\fbox{\centering
					\includegraphics[width=0.5\textwidth]{graphiken/Symboltable-inheritance.pdf}
				}
			\end{wrapfigure}
			
			Jedes Symbol hat einen Symboltypen. Die Symboltypen sind in der Klasse \texttt{SymbolType} definiert.
			Die vier bläulich hintergelben Interfaces in der Abbildung 4 haben keinen äquivalenten Symboltyp -- sie sind
			reine Metaklassen, die nur ein gemeinsames Interface der erbenden Klassen zusammenfassen, wobei
			es jedoch keine tatsächlichen Instanzen dieser Interfaces geben kann. Alle in der Abbildung grün umrandeten
			Interfaces erben von \texttt{QualifiedSymbol}.
			
			Das Interface der Symboltabelle ist streng hierarchisch aufgebaut. Es gibt die eine Basisklasse, von der
			alle anderen Interfaces (transitiv) erben. Diese strenge Hierarchie ist jedoch nur im Interface
			garantiert. In tatsächlichen Implementierungen könnten verschiedene Klassen zusammengefasst werden, um
			die Programmierung zu vereinfachen, weshalb eine Abfrage, ob ein Objekt Instanz einer bestimmten Symbolart
			nicht über \texttt{instanceof} erfolgen darf. Aus diesem Grund habe ich die Abfragemethode
			\verb|Boolean hasType(SymbolType leftType)| eingeführt. Diese gibt das Ergebnis in drei"=wertiger Logik
			zurück, wobei der Rückgabewert \texttt{null} gleichbedeutend mit unbekannt ist. Das Ergebnis \texttt{null}
			wird einerseits zurückgegeben, wenn das abgefragte Objekt unqualified ist oder \texttt{leftType} den Wert
			\texttt{UNQUALIFIED} trägt.
			
			Die Bedeutung des Parameters \texttt{leftType} ist hierbei so zu verstehen: Das abgefragte Symbol ist der
			\texttt{rightType}. Das Methode \texttt{hasType} wiederum entspricht der Implikation. Es wird also
			\texttt{true} zurückgegeben, falls $leftType\Rightarrow rightType$. Dies ist der Fall, wenn \texttt{leftType}
			in der Verberbungshierarchie über \texttt{rightType} steht, beziehungsweise, wenn beide Typen gleich sind.
			
			Die drei gelblich hinterlegten Interfaces sind Factories und können weitere Elemente erstellen, die
			sie dann enthalten. Außer durch die Factories können keine Objekte instanziiert werden. Instanzen können
			nicht an andere Container abgegeben werden. Dies soll versehentlichen Anwendungsfehlern vorbeugen.
			\begin{itemize}
				\item \texttt{Runtime}s können \texttt{Package}s erzeugen.
				\item \texttt{Package}s können \texttt{Class}es und \texttt{Interface}s erzeugen.
				\item \texttt{Class}es können \texttt{Member}s erzeugen.
				\item \texttt{Class}es und \texttt{Interface}s können \texttt{Method}s erzeugen, wobei jede Methode einen äußersten
					\texttt{Scope} enthält, abrufbar mit \texttt{Scope getScope()}.
				\item \texttt{Scope}s können \texttt{Variable}s erzeugen.
			\end{itemize}
			Sämtliche anderen Instanzen werden durch die \texttt{Runtime} erzeugt.
			
			Abfragen nach Objekten geschehen in \texttt{SymbolContainer}n und ihren Unterinterfaces. Zu diesem Zweck
			stellt das Interface die Methode
			\begin{itemize}
				\item \verb|Symbol tryGetQualifiedSymbol(PositionString name)|
			\end{itemize}
			bereit. Das zurückgegebene Symbol könnte qualified sein, falls das referenzierte Symbol bereits bekannt ist.
			Sollte das Symbol zum Zeitpunkt des Aufrufes noch nicht bekannt sein, so wird ein \texttt{UnqualifiedSymbol}
			zurückgegeben. Zu einem späteren Zeitpunkt kann ein unqualified Symbol mittels \texttt{QualifiedSymbol qualify()}
			qualifiziert werden. Sollte das Symbol dann immer noch nicht bekannt sein, so wird \texttt{null} zurückgegeben.
			Der Aufrufer muss in diesem Fall geeignet reagieren.
	
	\section{Anhang}
		\subsection{Abbildungs- und Quelltextverzeichnis}
			\listoffigures
			\lstlistoflistings
		
		\subsection{Lizenzierung}
			Teile des Quelltextes des Softwareprojektes Übersetzerbau 2010 stehen unter der Freien
			Lizenz \href{http://www.gnu.org/licenses/agpl-3.0.html}{GNU AGPL} in der Version 3.0.
			Die Dokumentation selbst, sowie die enthaltenen Graphiken, werden unter den Bedingungen
			der \href{http://www.gnu.org/licenses/fdl-1.3.html}{GNU FDL} in der Version 1.3 weitergegeben.
			Bindend sind nur die englischen Originaltext, die im nächsten Absatz angehängt sind.
			
			Sollten Sie den Quelltext zu anderen Konditionen weiterverwenden wollen, so setzen Sie
			sich bitte mit mir in Verbindung:
			\href{mailto:rene.kijewski@fu-berlin.de?subject=[SWP_CP_2010]}{rene.kijewski@fu"=berlin.de}.
			
			\subsubsection{Quelltext}
				\begin{multicols}{3}
					{\tiny\begin{center}
{\parindent 0in

Copyright \copyright\  2007 Free Software Foundation, Inc. \texttt{http://fsf.org/}

\bigskip
Everyone is permitted to copy and distribute verbatim copies of this

license document, but changing it is not allowed.}

\end{center}

\renewcommand{\abstractname}{Preamble}
The GNU Affero General Public License is a free, copyleft license
for software and other kinds of works, specifically designed to ensure
cooperation with the community in the case of network server software.

The licenses for most software and other practical works are
designed to take away your freedom to share and change the works.  By
contrast, our General Public Licenses are intended to guarantee your
freedom to share and change all versions of a program--to make sure it
remains free software for all its users.

When we speak of free software, we are referring to freedom, not
price.  Our General Public Licenses are designed to make sure that you
have the freedom to distribute copies of free software (and charge for
them if you wish), that you receive source code or can get it if you
want it, that you can change the software or use pieces of it in new
free programs, and that you know you can do these things.

Developers that use our General Public Licenses protect your rights
with two steps: (1) assert copyright on the software, and (2) offer
you this License which gives you legal permission to copy, distribute
and/or modify the software.

A secondary benefit of defending all users' freedom is that
improvements made in alternate versions of the program, if they
receive widespread use, become available for other developers to
incorporate.  Many developers of free software are heartened and
encouraged by the resulting cooperation.  However, in the case of
software used on network servers, this result may fail to come about.
The GNU General Public License permits making a modified version and
letting the public access it on a server without ever releasing its
source code to the public.

The GNU Affero General Public License is designed specifically to
ensure that, in such cases, the modified source code becomes available
to the community.  It requires the operator of a network server to
provide the source code of the modified version running there to the
users of that server.  Therefore, public use of a modified version, on
a publicly accessible server, gives the public access to the source
code of the modified version.

An older license, called the Affero General Public License and
published by Affero, was designed to accomplish similar goals.  This is
a different license, not a version of the Affero GPL, but Affero has
released a new version of the Affero GPL which permits relicensing under
this license.

The precise terms and conditions for copying, distribution and
modification follow.

\begin{center}
{\sc Terms and Conditions}
\end{center}

\begin{enumerate}

\addtocounter{enumi}{-1}

\item Definitions.

``This License'' refers to version 3 of the GNU Affero General Public License.

``Copyright'' also means copyright-like laws that apply to other kinds of
works, such as semiconductor masks.

``The Program'' refers to any copyrightable work licensed under this
License.  Each licensee is addressed as ``you''.  ``Licensees'' and
``recipients'' may be individuals or organizations.

To ``modify'' a work means to copy from or adapt all or part of the work
in a fashion requiring copyright permission, other than the making of an
exact copy.  The resulting work is called a ``modified version'' of the
earlier work or a work ``based on'' the earlier work.

A ``covered work'' means either the unmodified Program or a work based
on the Program.

To ``propagate'' a work means to do anything with it that, without
permission, would make you directly or secondarily liable for
infringement under applicable copyright law, except executing it on a
computer or modifying a private copy.  Propagation includes copying,
distribution (with or without modification), making available to the
public, and in some countries other activities as well.

To ``convey'' a work means any kind of propagation that enables other
parties to make or receive copies.  Mere interaction with a user through
a computer network, with no transfer of a copy, is not conveying.

An interactive user interface displays ``Appropriate Legal Notices''
to the extent that it includes a convenient and prominently visible
feature that (1) displays an appropriate copyright notice, and (2)
tells the user that there is no warranty for the work (except to the
extent that warranties are provided), that licensees may convey the
work under this License, and how to view a copy of this License.  If
the interface presents a list of user commands or options, such as a
menu, a prominent item in the list meets this criterion.

\item Source Code.

The ``source code'' for a work means the preferred form of the work
for making modifications to it.  ``Object code'' means any non-source
form of a work.

A ``Standard Interface'' means an interface that either is an official
standard defined by a recognized standards body, or, in the case of
interfaces specified for a particular programming language, one that
is widely used among developers working in that language.

The ``System Libraries'' of an executable work include anything, other
than the work as a whole, that (a) is included in the normal form of
packaging a Major Component, but which is not part of that Major
Component, and (b) serves only to enable use of the work with that
Major Component, or to implement a Standard Interface for which an
implementation is available to the public in source code form.  A
``Major Component'', in this context, means a major essential component
(kernel, window system, and so on) of the specific operating system
(if any) on which the executable work runs, or a compiler used to
produce the work, or an object code interpreter used to run it.

The ``Corresponding Source'' for a work in object code form means all
the source code needed to generate, install, and (for an executable
work) run the object code and to modify the work, including scripts to
control those activities.  However, it does not include the work's
System Libraries, or general-purpose tools or generally available free
programs which are used unmodified in performing those activities but
which are not part of the work.  For example, Corresponding Source
includes interface definition files associated with source files for
the work, and the source code for shared libraries and dynamically
linked subprograms that the work is specifically designed to require,
such as by intimate data communication or control flow between those
subprograms and other parts of the work.

The Corresponding Source need not include anything that users
can regenerate automatically from other parts of the Corresponding
Source.

The Corresponding Source for a work in source code form is that
same work.

\item Basic Permissions.

All rights granted under this License are granted for the term of
copyright on the Program, and are irrevocable provided the stated
conditions are met.  This License explicitly affirms your unlimited
permission to run the unmodified Program.  The output from running a
covered work is covered by this License only if the output, given its
content, constitutes a covered work.  This License acknowledges your
rights of fair use or other equivalent, as provided by copyright law.

You may make, run and propagate covered works that you do not
convey, without conditions so long as your license otherwise remains
in force.  You may convey covered works to others for the sole purpose
of having them make modifications exclusively for you, or provide you
with facilities for running those works, provided that you comply with
the terms of this License in conveying all material for which you do
not control copyright.  Those thus making or running the covered works
for you must do so exclusively on your behalf, under your direction
and control, on terms that prohibit them from making any copies of
your copyrighted material outside their relationship with you.

Conveying under any other circumstances is permitted solely under
the conditions stated below.  Sublicensing is not allowed; section 10
makes it unnecessary.

\item Protecting Users' Legal Rights From Anti-Circumvention Law.

No covered work shall be deemed part of an effective technological
measure under any applicable law fulfilling obligations under article
11 of the WIPO copyright treaty adopted on 20 December 1996, or
similar laws prohibiting or restricting circumvention of such
measures.

When you convey a covered work, you waive any legal power to forbid
circumvention of technological measures to the extent such circumvention
is effected by exercising rights under this License with respect to
the covered work, and you disclaim any intention to limit operation or
modification of the work as a means of enforcing, against the work's
users, your or third parties' legal rights to forbid circumvention of
technological measures.

\item Conveying Verbatim Copies.

You may convey verbatim copies of the Program's source code as you
receive it, in any medium, provided that you conspicuously and
appropriately publish on each copy an appropriate copyright notice;
keep intact all notices stating that this License and any
non-permissive terms added in accord with section 7 apply to the code;
keep intact all notices of the absence of any warranty; and give all
recipients a copy of this License along with the Program.

You may charge any price or no price for each copy that you convey,
and you may offer support or warranty protection for a fee.

\item Conveying Modified Source Versions.

You may convey a work based on the Program, or the modifications to
produce it from the Program, in the form of source code under the
terms of section 4, provided that you also meet all of these conditions:
  \begin{enumerate}
  \item The work must carry prominent notices stating that you modified
  it, and giving a relevant date.

  \item The work must carry prominent notices stating that it is
  released under this License and any conditions added under section
  7.  This requirement modifies the requirement in section 4 to
  ``keep intact all notices''.

  \item You must license the entire work, as a whole, under this
  License to anyone who comes into possession of a copy.  This
  License will therefore apply, along with any applicable section 7
  additional terms, to the whole of the work, and all its parts,
  regardless of how they are packaged.  This License gives no
  permission to license the work in any other way, but it does not
  invalidate such permission if you have separately received it.

  \item If the work has interactive user interfaces, each must display
  Appropriate Legal Notices; however, if the Program has interactive
  interfaces that do not display Appropriate Legal Notices, your
  work need not make them do so.
\end{enumerate}
A compilation of a covered work with other separate and independent
works, which are not by their nature extensions of the covered work,
and which are not combined with it such as to form a larger program,
in or on a volume of a storage or distribution medium, is called an
``aggregate'' if the compilation and its resulting copyright are not
used to limit the access or legal rights of the compilation's users
beyond what the individual works permit.  Inclusion of a covered work
in an aggregate does not cause this License to apply to the other
parts of the aggregate.

\item Conveying Non-Source Forms.

You may convey a covered work in object code form under the terms
of sections 4 and 5, provided that you also convey the
machine-readable Corresponding Source under the terms of this License,
in one of these ways:
  \begin{enumerate}
  \item Convey the object code in, or embodied in, a physical product
  (including a physical distribution medium), accompanied by the
  Corresponding Source fixed on a durable physical medium
  customarily used for software interchange.

  \item Convey the object code in, or embodied in, a physical product
  (including a physical distribution medium), accompanied by a
  written offer, valid for at least three years and valid for as
  long as you offer spare parts or customer support for that product
  model, to give anyone who possesses the object code either (1) a
  copy of the Corresponding Source for all the software in the
  product that is covered by this License, on a durable physical
  medium customarily used for software interchange, for a price no
  more than your reasonable cost of physically performing this
  conveying of source, or (2) access to copy the
  Corresponding Source from a network server at no charge.

  \item Convey individual copies of the object code with a copy of the
  written offer to provide the Corresponding Source.  This
  alternative is allowed only occasionally and noncommercially, and
  only if you received the object code with such an offer, in accord
  with subsection 6b.

  \item Convey the object code by offering access from a designated
  place (gratis or for a charge), and offer equivalent access to the
  Corresponding Source in the same way through the same place at no
  further charge.  You need not require recipients to copy the
  Corresponding Source along with the object code.  If the place to
  copy the object code is a network server, the Corresponding Source
  may be on a different server (operated by you or a third party)
  that supports equivalent copying facilities, provided you maintain
  clear directions next to the object code saying where to find the
  Corresponding Source.  Regardless of what server hosts the
  Corresponding Source, you remain obligated to ensure that it is
  available for as long as needed to satisfy these requirements.

  \item Convey the object code using peer-to-peer transmission, provided
  you inform other peers where the object code and Corresponding
  Source of the work are being offered to the general public at no
  charge under subsection 6d.
  \end{enumerate}

A separable portion of the object code, whose source code is excluded
from the Corresponding Source as a System Library, need not be
included in conveying the object code work.

A ``User Product'' is either (1) a ``consumer product'', which means any
tangible personal property which is normally used for personal, family,
or household purposes, or (2) anything designed or sold for incorporation
into a dwelling.  In determining whether a product is a consumer product,
doubtful cases shall be resolved in favor of coverage.  For a particular
product received by a particular user, ``normally used'' refers to a
typical or common use of that class of product, regardless of the status
of the particular user or of the way in which the particular user
actually uses, or expects or is expected to use, the product.  A product
is a consumer product regardless of whether the product has substantial
commercial, industrial or non-consumer uses, unless such uses represent
the only significant mode of use of the product.

``Installation Information'' for a User Product means any methods,
procedures, authorization keys, or other information required to install
and execute modified versions of a covered work in that User Product from
a modified version of its Corresponding Source.  The information must
suffice to ensure that the continued functioning of the modified object
code is in no case prevented or interfered with solely because
modification has been made.

If you convey an object code work under this section in, or with, or
specifically for use in, a User Product, and the conveying occurs as
part of a transaction in which the right of possession and use of the
User Product is transferred to the recipient in perpetuity or for a
fixed term (regardless of how the transaction is characterized), the
Corresponding Source conveyed under this section must be accompanied
by the Installation Information.  But this requirement does not apply
if neither you nor any third party retains the ability to install
modified object code on the User Product (for example, the work has
been installed in ROM).

The requirement to provide Installation Information does not include a
requirement to continue to provide support service, warranty, or updates
for a work that has been modified or installed by the recipient, or for
the User Product in which it has been modified or installed.  Access to a
network may be denied when the modification itself materially and
adversely affects the operation of the network or violates the rules and
protocols for communication across the network.

Corresponding Source conveyed, and Installation Information provided,
in accord with this section must be in a format that is publicly
documented (and with an implementation available to the public in
source code form), and must require no special password or key for
unpacking, reading or copying.

\item Additional Terms.

``Additional permissions'' are terms that supplement the terms of this
License by making exceptions from one or more of its conditions.
Additional permissions that are applicable to the entire Program shall
be treated as though they were included in this License, to the extent
that they are valid under applicable law.  If additional permissions
apply only to part of the Program, that part may be used separately
under those permissions, but the entire Program remains governed by
this License without regard to the additional permissions.

When you convey a copy of a covered work, you may at your option
remove any additional permissions from that copy, or from any part of
it.  (Additional permissions may be written to require their own
removal in certain cases when you modify the work.)  You may place
additional permissions on material, added by you to a covered work,
for which you have or can give appropriate copyright permission.

Notwithstanding any other provision of this License, for material you
add to a covered work, you may (if authorized by the copyright holders of
that material) supplement the terms of this License with terms:
  \begin{enumerate}
  \item Disclaiming warranty or limiting liability differently from the
  terms of sections 15 and 16 of this License; or

  \item Requiring preservation of specified reasonable legal notices or
  author attributions in that material or in the Appropriate Legal
  Notices displayed by works containing it; or

  \item Prohibiting misrepresentation of the origin of that material, or
  requiring that modified versions of such material be marked in
  reasonable ways as different from the original version; or

  \item Limiting the use for publicity purposes of names of licensors or
  authors of the material; or

  \item Declining to grant rights under trademark law for use of some
  trade names, trademarks, or service marks; or

  \item Requiring indemnification of licensors and authors of that
  material by anyone who conveys the material (or modified versions of
  it) with contractual assumptions of liability to the recipient, for
  any liability that these contractual assumptions directly impose on
  those licensors and authors.
  \end{enumerate}

All other non-permissive additional terms are considered ``further
restrictions'' within the meaning of section 10.  If the Program as you
received it, or any part of it, contains a notice stating that it is
governed by this License along with a term that is a further
restriction, you may remove that term.  If a license document contains
a further restriction but permits relicensing or conveying under this
License, you may add to a covered work material governed by the terms
of that license document, provided that the further restriction does
not survive such relicensing or conveying.

If you add terms to a covered work in accord with this section, you
must place, in the relevant source files, a statement of the
additional terms that apply to those files, or a notice indicating
where to find the applicable terms.

Additional terms, permissive or non-permissive, may be stated in the
form of a separately written license, or stated as exceptions;
the above requirements apply either way.

\item Termination.

You may not propagate or modify a covered work except as expressly
provided under this License.  Any attempt otherwise to propagate or
modify it is void, and will automatically terminate your rights under
this License (including any patent licenses granted under the third
paragraph of section 11).

However, if you cease all violation of this License, then your
license from a particular copyright holder is reinstated (a)
provisionally, unless and until the copyright holder explicitly and
finally terminates your license, and (b) permanently, if the copyright
holder fails to notify you of the violation by some reasonable means
prior to 60 days after the cessation.

Moreover, your license from a particular copyright holder is
reinstated permanently if the copyright holder notifies you of the
violation by some reasonable means, this is the first time you have
received notice of violation of this License (for any work) from that
copyright holder, and you cure the violation prior to 30 days after
your receipt of the notice.

Termination of your rights under this section does not terminate the
licenses of parties who have received copies or rights from you under
this License.  If your rights have been terminated and not permanently
reinstated, you do not qualify to receive new licenses for the same
material under section 10.

\item Acceptance Not Required for Having Copies.

You are not required to accept this License in order to receive or
run a copy of the Program.  Ancillary propagation of a covered work
occurring solely as a consequence of using peer-to-peer transmission
to receive a copy likewise does not require acceptance.  However,
nothing other than this License grants you permission to propagate or
modify any covered work.  These actions infringe copyright if you do
not accept this License.  Therefore, by modifying or propagating a
covered work, you indicate your acceptance of this License to do so.

\item Automatic Licensing of Downstream Recipients.

Each time you convey a covered work, the recipient automatically
receives a license from the original licensors, to run, modify and
propagate that work, subject to this License.  You are not responsible
for enforcing compliance by third parties with this License.

An ``entity transaction'' is a transaction transferring control of an
organization, or substantially all assets of one, or subdividing an
organization, or merging organizations.  If propagation of a covered
work results from an entity transaction, each party to that
transaction who receives a copy of the work also receives whatever
licenses to the work the party's predecessor in interest had or could
give under the previous paragraph, plus a right to possession of the
Corresponding Source of the work from the predecessor in interest, if
the predecessor has it or can get it with reasonable efforts.

You may not impose any further restrictions on the exercise of the
rights granted or affirmed under this License.  For example, you may
not impose a license fee, royalty, or other charge for exercise of
rights granted under this License, and you may not initiate litigation
(including a cross-claim or counterclaim in a lawsuit) alleging that
any patent claim is infringed by making, using, selling, offering for
sale, or importing the Program or any portion of it.

\item Patents.

A ``contributor'' is a copyright holder who authorizes use under this
License of the Program or a work on which the Program is based.  The
work thus licensed is called the contributor's ``contributor version''.

A contributor's ``essential patent claims'' are all patent claims
owned or controlled by the contributor, whether already acquired or
hereafter acquired, that would be infringed by some manner, permitted
by this License, of making, using, or selling its contributor version,
but do not include claims that would be infringed only as a
consequence of further modification of the contributor version.  For
purposes of this definition, ``control'' includes the right to grant
patent sublicenses in a manner consistent with the requirements of
this License.

Each contributor grants you a non-exclusive, worldwide, royalty-free
patent license under the contributor's essential patent claims, to
make, use, sell, offer for sale, import and otherwise run, modify and
propagate the contents of its contributor version.

In the following three paragraphs, a ``patent license'' is any express
agreement or commitment, however denominated, not to enforce a patent
(such as an express permission to practice a patent or covenant not to
sue for patent infringement).  To ``grant'' such a patent license to a
party means to make such an agreement or commitment not to enforce a
patent against the party.

If you convey a covered work, knowingly relying on a patent license,
and the Corresponding Source of the work is not available for anyone
to copy, free of charge and under the terms of this License, through a
publicly available network server or other readily accessible means,
then you must either (1) cause the Corresponding Source to be so
available, or (2) arrange to deprive yourself of the benefit of the
patent license for this particular work, or (3) arrange, in a manner
consistent with the requirements of this License, to extend the patent
license to downstream recipients.  ``Knowingly relying'' means you have
actual knowledge that, but for the patent license, your conveying the
covered work in a country, or your recipient's use of the covered work
in a country, would infringe one or more identifiable patents in that
country that you have reason to believe are valid.

If, pursuant to or in connection with a single transaction or
arrangement, you convey, or propagate by procuring conveyance of, a
covered work, and grant a patent license to some of the parties
receiving the covered work authorizing them to use, propagate, modify
or convey a specific copy of the covered work, then the patent license
you grant is automatically extended to all recipients of the covered
work and works based on it.

A patent license is ``discriminatory'' if it does not include within
the scope of its coverage, prohibits the exercise of, or is
conditioned on the non-exercise of one or more of the rights that are
specifically granted under this License.  You may not convey a covered
work if you are a party to an arrangement with a third party that is
in the business of distributing software, under which you make payment
to the third party based on the extent of your activity of conveying
the work, and under which the third party grants, to any of the
parties who would receive the covered work from you, a discriminatory
patent license (a) in connection with copies of the covered work
conveyed by you (or copies made from those copies), or (b) primarily
for and in connection with specific products or compilations that
contain the covered work, unless you entered into that arrangement,
or that patent license was granted, prior to 28 March 2007.

Nothing in this License shall be construed as excluding or limiting
any implied license or other defenses to infringement that may
otherwise be available to you under applicable patent law.

\item No Surrender of Others' Freedom.

If conditions are imposed on you (whether by court order, agreement or
otherwise) that contradict the conditions of this License, they do not
excuse you from the conditions of this License.  If you cannot convey a
covered work so as to satisfy simultaneously your obligations under this
License and any other pertinent obligations, then as a consequence you may
not convey it at all.  For example, if you agree to terms that obligate you
to collect a royalty for further conveying from those to whom you convey
the Program, the only way you could satisfy both those terms and this
License would be to refrain entirely from conveying the Program.

\item Remote Network Interaction; Use with the GNU General Public License.

Notwithstanding any other provision of this License, if you modify the
Program, your modified version must prominently offer all users interacting
with it remotely through a computer network (if your version supports such
interaction) an opportunity to receive the Corresponding Source of your
version by providing access to the Corresponding Source from a network
server at no charge, through some standard or customary means of
facilitating copying of software.  This Corresponding Source shall include
the Corresponding Source for any work covered by version 3 of the GNU
General Public License that is incorporated pursuant to the following
paragraph.

Notwithstanding any other provision of this License, you have permission to
link or combine any covered work with a work licensed under version 3 of
the GNU General Public License into a single combined work, and to convey
the resulting work.  The terms of this License will continue to apply to
the part which is the covered work, but the work with which it is combined
will remain governed by version 3 of the GNU General Public License.

\item Revised Versions of this License.

The Free Software Foundation may publish revised and/or new versions of
the GNU Affero General Public License from time to time.  Such new versions will
be similar in spirit to the present version, but may differ in detail to
address new problems or concerns.

Each version is given a distinguishing version number.  If the
Program specifies that a certain numbered version of the GNU Affero General
Public License ``or any later version'' applies to it, you have the
option of following the terms and conditions either of that numbered
version or of any later version published by the Free Software
Foundation.  If the Program does not specify a version number of the
GNU Affero General Public License, you may choose any version ever published
by the Free Software Foundation.

If the Program specifies that a proxy can decide which future
versions of the GNU Affero General Public License can be used, that proxy's
public statement of acceptance of a version permanently authorizes you
to choose that version for the Program.

Later license versions may give you additional or different
permissions.  However, no additional obligations are imposed on any
author or copyright holder as a result of your choosing to follow a
later version.

\item Disclaimer of Warranty.

\begin{sloppypar}
 THERE IS NO WARRANTY FOR THE PROGRAM, TO THE EXTENT PERMITTED BY
 APPLICABLE LAW.  EXCEPT WHEN OTHERWISE STATED IN WRITING THE
 COPYRIGHT HOLDERS AND/OR OTHER PARTIES PROVIDE THE PROGRAM ``AS IS''
 WITHOUT WARRANTY OF ANY KIND, EITHER EXPRESSED OR IMPLIED,
 INCLUDING, BUT NOT LIMITED TO, THE IMPLIED WARRANTIES OF
 MERCHANTABILITY AND FITNESS FOR A PARTICULAR PURPOSE.  THE ENTIRE
 RISK AS TO THE QUALITY AND PERFORMANCE OF THE PROGRAM IS WITH YOU.
 SHOULD THE PROGRAM PROVE DEFECTIVE, YOU ASSUME THE COST OF ALL
 NECESSARY SERVICING, REPAIR OR CORRECTION.
\end{sloppypar}

\item Limitation of Liability.

 IN NO EVENT UNLESS REQUIRED BY APPLICABLE LAW OR AGREED TO IN
 WRITING WILL ANY COPYRIGHT HOLDER, OR ANY OTHER PARTY WHO MODIFIES
 AND/OR CONVEYS THE PROGRAM AS PERMITTED ABOVE, BE LIABLE TO YOU FOR
 DAMAGES, INCLUDING ANY GENERAL, SPECIAL, INCIDENTAL OR CONSEQUENTIAL
 DAMAGES ARISING OUT OF THE USE OR INABILITY TO USE THE PROGRAM
 (INCLUDING BUT NOT LIMITED TO LOSS OF DATA OR DATA BEING RENDERED
 INACCURATE OR LOSSES SUSTAINED BY YOU OR THIRD PARTIES OR A FAILURE
 OF THE PROGRAM TO OPERATE WITH ANY OTHER PROGRAMS), EVEN IF SUCH
 HOLDER OR OTHER PARTY HAS BEEN ADVISED OF THE POSSIBILITY OF SUCH
 DAMAGES.

\item Interpretation of Sections 15 and 16.

If the disclaimer of warranty and limitation of liability provided
above cannot be given local legal effect according to their terms,
reviewing courts shall apply local law that most closely approximates
an absolute waiver of all civil liability in connection with the
Program, unless a warranty or assumption of liability accompanies a
copy of the Program in return for a fee.

\begin{center}
{\sc End of Terms and Conditions}

\bigskip
How to Apply These Terms to Your New Programs
\end{center}

If you develop a new program, and you want it to be of the greatest
possible use to the public, the best way to achieve this is to make it
free software which everyone can redistribute and change under these terms.

To do so, attach the following notices to the program.  It is safest
to attach them to the start of each source file to most effectively
state the exclusion of warranty; and each file should have at least
the ``copyright'' line and a pointer to where the full notice is found.

<one line to give the program's name and a brief idea of what it does.>

Copyright (C) <textyear>  <name of author>

This program is free software: you can redistribute it and/or modify
it under the terms of the GNU Affero General Public License as published by
the Free Software Foundation, either version 3 of the License, or
(at your option) any later version.

This program is distributed in the hope that it will be useful,
but WITHOUT ANY WARRANTY; without even the implied warranty of
MERCHANTABILITY or FITNESS FOR A PARTICULAR PURPOSE.  See the
GNU Affero General Public License for more details.

You should have received a copy of the GNU Affero General Public License
along with this program.  If not, see <http://www.gnu.org/licenses/>.

Also add information on how to contact you by electronic and paper mail.

If your software can interact with users remotely through a computer
network, you should also make sure that it provides a way for users to
get its source.  For example, if your program is a web application, its
interface could display a ``Source'' link that leads users to an archive
of the code.  There are many ways you could offer source, and different
solutions will be better for different programs; see section 13 for the
specific requirements.

You should also get your employer (if you work as a programmer) or
school, if any, to sign a ``copyright disclaimer'' for the program, if
necessary.  For more information on this, and how to apply and follow
the GNU AGPL, see \texttt{http://www.gnu.org/licenses/}.

\end{enumerate}
}
				\end{multicols}
			
			\subsubsection[Dieses Dokument]{Dieses Dokument}
				\begin{multicols}{3}
					{\tiny\begin{center}
		 Version 1.3, 3 November 2008

 Copyright \copyright{} 2000, 2001, 2002, 2007, 2008  Free Software Foundation, Inc.
 
 \bigskip
 
	  <http://fsf.org/>
  
 \bigskip
 
 Everyone is permitted to copy and distribute verbatim copies
 of this license document, but changing it is not allowed.
\end{center}


\begin{center}
{\bf Preamble}
\end{center}

The purpose of this License is to make a manual, textbook, or other
functional and useful document ``free'' in the sense of freedom: to
assure everyone the effective freedom to copy and redistribute it,
with or without modifying it, either commercially or noncommercially.
Secondarily, this License preserves for the author and publisher a way
to get credit for their work, while not being considered responsible
for modifications made by others.

This License is a kind of ``copyleft'', which means that derivative
works of the document must themselves be free in the same sense.  It
complements the GNU General Public License, which is a copyleft
license designed for free software.

We have designed this License in order to use it for manuals for free
software, because free software needs free documentation: a free
program should come with manuals providing the same freedoms that the
software does.  But this License is not limited to software manuals;
it can be used for any textual work, regardless of subject matter or
whether it is published as a printed book.  We recommend this License
principally for works whose purpose is instruction or reference.


\begin{center}
{\bf 1. APPLICABILITY AND DEFINITIONS\par}
\phantomsection
\end{center}

This License applies to any manual or other work, in any medium, that
contains a notice placed by the copyright holder saying it can be
distributed under the terms of this License.  Such a notice grants a
world-wide, royalty-free license, unlimited in duration, to use that
work under the conditions stated herein.  The ``\textbf{Document}'', below,
refers to any such manual or work.  Any member of the public is a
licensee, and is addressed as ``\textbf{you}''.  You accept the license if you
copy, modify or distribute the work in a way requiring permission
under copyright law.

A ``\textbf{Modified Version}'' of the Document means any work containing the
Document or a portion of it, either copied verbatim, or with
modifications and/or translated into another language.

A ``\textbf{Secondary Section}'' is a named appendix or a front-matter section of
the Document that deals exclusively with the relationship of the
publishers or authors of the Document to the Document's overall subject
(or to related matters) and contains nothing that could fall directly
within that overall subject.  (Thus, if the Document is in part a
textbook of mathematics, a Secondary Section may not explain any
mathematics.)  The relationship could be a matter of historical
connection with the subject or with related matters, or of legal,
commercial, philosophical, ethical or political position regarding
them.

The ``\textbf{Invariant Sections}'' are certain Secondary Sections whose titles
are designated, as being those of Invariant Sections, in the notice
that says that the Document is released under this License.  If a
section does not fit the above definition of Secondary then it is not
allowed to be designated as Invariant.  The Document may contain zero
Invariant Sections.  If the Document does not identify any Invariant
Sections then there are none.

The ``\textbf{Cover Texts}'' are certain short passages of text that are listed,
as Front-Cover Texts or Back-Cover Texts, in the notice that says that
the Document is released under this License.  A Front-Cover Text may
be at most 5 words, and a Back-Cover Text may be at most 25 words.

A ``\textbf{Transparent}'' copy of the Document means a machine-readable copy,
represented in a format whose specification is available to the
general public, that is suitable for revising the document
straightforwardly with generic text editors or (for images composed of
pixels) generic paint programs or (for drawings) some widely available
drawing editor, and that is suitable for input to text formatters or
for automatic translation to a variety of formats suitable for input
to text formatters.  A copy made in an otherwise Transparent file
format whose markup, or absence of markup, has been arranged to thwart
or discourage subsequent modification by readers is not Transparent.
An image format is not Transparent if used for any substantial amount
of text.  A copy that is not ``Transparent'' is called ``\textbf{Opaque}''.

Examples of suitable formats for Transparent copies include plain
ASCII without markup, Texinfo input format, LaTeX input format, SGML
or XML using a publicly available DTD, and standard-conforming simple
HTML, PostScript or PDF designed for human modification.  Examples of
transparent image formats include PNG, XCF and JPG.  Opaque formats
include proprietary formats that can be read and edited only by
proprietary word processors, SGML or XML for which the DTD and/or
processing tools are not generally available, and the
machine-generated HTML, PostScript or PDF produced by some word
processors for output purposes only.

The ``\textbf{Title Page}'' means, for a printed book, the title page itself,
plus such following pages as are needed to hold, legibly, the material
this License requires to appear in the title page.  For works in
formats which do not have any title page as such, ``Title Page'' means
the text near the most prominent appearance of the work's title,
preceding the beginning of the body of the text.

The ``\textbf{publisher}'' means any person or entity that distributes
copies of the Document to the public.

A section ``\textbf{Entitled XYZ}'' means a named subunit of the Document whose
title either is precisely XYZ or contains XYZ in parentheses following
text that translates XYZ in another language.  (Here XYZ stands for a
specific section name mentioned below, such as ``\textbf{Acknowledgements}'',
``\textbf{Dedications}'', ``\textbf{Endorsements}'', or ``\textbf{History}''.)  
To ``\textbf{Preserve the Title}''
of such a section when you modify the Document means that it remains a
section ``Entitled XYZ'' according to this definition.

The Document may include Warranty Disclaimers next to the notice which
states that this License applies to the Document.  These Warranty
Disclaimers are considered to be included by reference in this
License, but only as regards disclaiming warranties: any other
implication that these Warranty Disclaimers may have is void and has
no effect on the meaning of this License.


\begin{center}
{\bf 2. VERBATIM COPYING\par}
\phantomsection
\end{center}

You may copy and distribute the Document in any medium, either
commercially or noncommercially, provided that this License, the
copyright notices, and the license notice saying this License applies
to the Document are reproduced in all copies, and that you add no other
conditions whatsoever to those of this License.  You may not use
technical measures to obstruct or control the reading or further
copying of the copies you make or distribute.  However, you may accept
compensation in exchange for copies.  If you distribute a large enough
number of copies you must also follow the conditions in section~3.

You may also lend copies, under the same conditions stated above, and
you may publicly display copies.


\begin{center}
{\bf 3. COPYING IN QUANTITY\par}
\phantomsection
\end{center}

If you publish printed copies (or copies in media that commonly have
printed covers) of the Document, numbering more than 100, and the
Document's license notice requires Cover Texts, you must enclose the
copies in covers that carry, clearly and legibly, all these Cover
Texts: Front-Cover Texts on the front cover, and Back-Cover Texts on
the back cover.  Both covers must also clearly and legibly identify
you as the publisher of these copies.  The front cover must present
the full title with all words of the title equally prominent and
visible.  You may add other material on the covers in addition.
Copying with changes limited to the covers, as long as they preserve
the title of the Document and satisfy these conditions, can be treated
as verbatim copying in other respects.

If the required texts for either cover are too voluminous to fit
legibly, you should put the first ones listed (as many as fit
reasonably) on the actual cover, and continue the rest onto adjacent
pages.

If you publish or distribute Opaque copies of the Document numbering
more than 100, you must either include a machine-readable Transparent
copy along with each Opaque copy, or state in or with each Opaque copy
a computer-network location from which the general network-using
public has access to download using public-standard network protocols
a complete Transparent copy of the Document, free of added material.
If you use the latter option, you must take reasonably prudent steps,
when you begin distribution of Opaque copies in quantity, to ensure
that this Transparent copy will remain thus accessible at the stated
location until at least one year after the last time you distribute an
Opaque copy (directly or through your agents or retailers) of that
edition to the public.

It is requested, but not required, that you contact the authors of the
Document well before redistributing any large number of copies, to give
them a chance to provide you with an updated version of the Document.


\begin{center}
{\bf 4. MODIFICATIONS\par}
\phantomsection
\end{center}

You may copy and distribute a Modified Version of the Document under
the conditions of sections 2 and 3 above, provided that you release
the Modified Version under precisely this License, with the Modified
Version filling the role of the Document, thus licensing distribution
and modification of the Modified Version to whoever possesses a copy
of it.  In addition, you must do these things in the Modified Version:

\begin{itemize}
\item[A.] 
	Use in the Title Page (and on the covers, if any) a title distinct
	from that of the Document, and from those of previous versions
	(which should, if there were any, be listed in the History section
	of the Document).  You may use the same title as a previous version
	if the original publisher of that version gives permission.
	
\item[B.]
	List on the Title Page, as authors, one or more persons or entities
	responsible for authorship of the modifications in the Modified
	Version, together with at least five of the principal authors of the
	Document (all of its principal authors, if it has fewer than five),
	unless they release you from this requirement.
	
\item[C.]
	State on the Title page the name of the publisher of the
	Modified Version, as the publisher.
	
\item[D.]
	Preserve all the copyright notices of the Document.
	
\item[E.]
	Add an appropriate copyright notice for your modifications
	adjacent to the other copyright notices.
	
\item[F.]
	Include, immediately after the copyright notices, a license notice
	giving the public permission to use the Modified Version under the
	terms of this License, in the form shown in the Addendum below.
	
\item[G.]
	Preserve in that license notice the full lists of Invariant Sections
	and required Cover Texts given in the Document's license notice.
	
\item[H.]
	Include an unaltered copy of this License.
	
\item[I.]
	Preserve the section Entitled ``History'', Preserve its Title, and add
	to it an item stating at least the title, year, new authors, and
	publisher of the Modified Version as given on the Title Page.  If
	there is no section Entitled ``History'' in the Document, create one
	stating the title, year, authors, and publisher of the Document as
	given on its Title Page, then add an item describing the Modified
	Version as stated in the previous sentence.
	
\item[J.]
	Preserve the network location, if any, given in the Document for
	public access to a Transparent copy of the Document, and likewise
	the network locations given in the Document for previous versions
	it was based on.  These may be placed in the ``History'' section.
	You may omit a network location for a work that was published at
	least four years before the Document itself, or if the original
	publisher of the version it refers to gives permission.
	
\item[K.]
	For any section Entitled ``Acknowledgements'' or ``Dedications'',
	Preserve the Title of the section, and preserve in the section all
	the substance and tone of each of the contributor acknowledgements
	and/or dedications given therein.
	
\item[L.]
	Preserve all the Invariant Sections of the Document,
	unaltered in their text and in their titles.  Section numbers
	or the equivalent are not considered part of the section titles.
	
\item[M.]
	Delete any section Entitled ``Endorsements''.  Such a section
	may not be included in the Modified Version.
	
\item[N.]
	Do not retitle any existing section to be Entitled ``Endorsements''
	or to conflict in title with any Invariant Section.
	
\item[O.]
	Preserve any Warranty Disclaimers.
\end{itemize}

If the Modified Version includes new front-matter sections or
appendices that qualify as Secondary Sections and contain no material
copied from the Document, you may at your option designate some or all
of these sections as invariant.  To do this, add their titles to the
list of Invariant Sections in the Modified Version's license notice.
These titles must be distinct from any other section titles.

You may add a section Entitled ``Endorsements'', provided it contains
nothing but endorsements of your Modified Version by various
parties---for example, statements of peer review or that the text has
been approved by an organization as the authoritative definition of a
standard.

You may add a passage of up to five words as a Front-Cover Text, and a
passage of up to 25 words as a Back-Cover Text, to the end of the list
of Cover Texts in the Modified Version.  Only one passage of
Front-Cover Text and one of Back-Cover Text may be added by (or
through arrangements made by) any one entity.  If the Document already
includes a cover text for the same cover, previously added by you or
by arrangement made by the same entity you are acting on behalf of,
you may not add another; but you may replace the old one, on explicit
permission from the previous publisher that added the old one.

The author(s) and publisher(s) of the Document do not by this License
give permission to use their names for publicity for or to assert or
imply endorsement of any Modified Version.


\begin{center}
{\bf 5. COMBINING DOCUMENTS\par}
\phantomsection
\end{center}

You may combine the Document with other documents released under this
License, under the terms defined in section~4 above for modified
versions, provided that you include in the combination all of the
Invariant Sections of all of the original documents, unmodified, and
list them all as Invariant Sections of your combined work in its
license notice, and that you preserve all their Warranty Disclaimers.

The combined work need only contain one copy of this License, and
multiple identical Invariant Sections may be replaced with a single
copy.  If there are multiple Invariant Sections with the same name but
different contents, make the title of each such section unique by
adding at the end of it, in parentheses, the name of the original
author or publisher of that section if known, or else a unique number.
Make the same adjustment to the section titles in the list of
Invariant Sections in the license notice of the combined work.

In the combination, you must combine any sections Entitled ``History''
in the various original documents, forming one section Entitled
``History''; likewise combine any sections Entitled ``Acknowledgements'',
and any sections Entitled ``Dedications''.  You must delete all sections
Entitled ``Endorsements''.

\begin{center}
{\bf 6. COLLECTIONS OF DOCUMENTS\par}
\phantomsection
\end{center}

You may make a collection consisting of the Document and other documents
released under this License, and replace the individual copies of this
License in the various documents with a single copy that is included in
the collection, provided that you follow the rules of this License for
verbatim copying of each of the documents in all other respects.

You may extract a single document from such a collection, and distribute
it individually under this License, provided you insert a copy of this
License into the extracted document, and follow this License in all
other respects regarding verbatim copying of that document.


\begin{center}
{\bf 7. AGGREGATION WITH INDEPENDENT WORKS\par}
\phantomsection
\end{center}

A compilation of the Document or its derivatives with other separate
and independent documents or works, in or on a volume of a storage or
distribution medium, is called an ``aggregate'' if the copyright
resulting from the compilation is not used to limit the legal rights
of the compilation's users beyond what the individual works permit.
When the Document is included in an aggregate, this License does not
apply to the other works in the aggregate which are not themselves
derivative works of the Document.

If the Cover Text requirement of section~3 is applicable to these
copies of the Document, then if the Document is less than one half of
the entire aggregate, the Document's Cover Texts may be placed on
covers that bracket the Document within the aggregate, or the
electronic equivalent of covers if the Document is in electronic form.
Otherwise they must appear on printed covers that bracket the whole
aggregate.


\begin{center}
{\bf 8. TRANSLATION\par}
\phantomsection
\end{center}

Translation is considered a kind of modification, so you may
distribute translations of the Document under the terms of section~4.
Replacing Invariant Sections with translations requires special
permission from their copyright holders, but you may include
translations of some or all Invariant Sections in addition to the
original versions of these Invariant Sections.  You may include a
translation of this License, and all the license notices in the
Document, and any Warranty Disclaimers, provided that you also include
the original English version of this License and the original versions
of those notices and disclaimers.  In case of a disagreement between
the translation and the original version of this License or a notice
or disclaimer, the original version will prevail.

If a section in the Document is Entitled ``Acknowledgements'',
``Dedications'', or ``History'', the requirement (section~4) to Preserve
its Title (section~1) will typically require changing the actual
title.


\begin{center}
{\bf 9. TERMINATION\par}
\phantomsection
\end{center}

You may not copy, modify, sublicense, or distribute the Document
except as expressly provided under this License.  Any attempt
otherwise to copy, modify, sublicense, or distribute it is void, and
will automatically terminate your rights under this License.

However, if you cease all violation of this License, then your license
from a particular copyright holder is reinstated (a) provisionally,
unless and until the copyright holder explicitly and finally
terminates your license, and (b) permanently, if the copyright holder
fails to notify you of the violation by some reasonable means prior to
60 days after the cessation.

Moreover, your license from a particular copyright holder is
reinstated permanently if the copyright holder notifies you of the
violation by some reasonable means, this is the first time you have
received notice of violation of this License (for any work) from that
copyright holder, and you cure the violation prior to 30 days after
your receipt of the notice.

Termination of your rights under this section does not terminate the
licenses of parties who have received copies or rights from you under
this License.  If your rights have been terminated and not permanently
reinstated, receipt of a copy of some or all of the same material does
not give you any rights to use it.


\begin{center}
{\bf 10. FUTURE REVISIONS OF THIS LICENSE\par}
\phantomsection
\end{center}

The Free Software Foundation may publish new, revised versions
of the GNU Free Documentation License from time to time.  Such new
versions will be similar in spirit to the present version, but may
differ in detail to address new problems or concerns.  See
http://www.gnu.org/copyleft/.

Each version of the License is given a distinguishing version number.
If the Document specifies that a particular numbered version of this
License ``or any later version'' applies to it, you have the option of
following the terms and conditions either of that specified version or
of any later version that has been published (not as a draft) by the
Free Software Foundation.  If the Document does not specify a version
number of this License, you may choose any version ever published (not
as a draft) by the Free Software Foundation.  If the Document
specifies that a proxy can decide which future versions of this
License can be used, that proxy's public statement of acceptance of a
version permanently authorizes you to choose that version for the
Document.


\begin{center}
{\bf 11. RELICENSING\par}
\phantomsection
\end{center}

``Massive Multiauthor Collaboration Site'' (or ``MMC Site'') means any
World Wide Web server that publishes copyrightable works and also
provides prominent facilities for anybody to edit those works.  A
public wiki that anybody can edit is an example of such a server.  A
``Massive Multiauthor Collaboration'' (or ``MMC'') contained in the
site means any set of copyrightable works thus published on the MMC
site.

``CC-BY-SA'' means the Creative Commons Attribution-Share Alike 3.0
license published by Creative Commons Corporation, a not-for-profit
corporation with a principal place of business in San Francisco,
California, as well as future copyleft versions of that license
published by that same organization.

``Incorporate'' means to publish or republish a Document, in whole or
in part, as part of another Document.

An MMC is ``eligible for relicensing'' if it is licensed under this
License, and if all works that were first published under this License
somewhere other than this MMC, and subsequently incorporated in whole
or in part into the MMC, (1) had no cover texts or invariant sections,
and (2) were thus incorporated prior to November 1, 2008.

The operator of an MMC Site may republish an MMC contained in the site
under CC-BY-SA on the same site at any time before August 1, 2009,
provided the MMC is eligible for relicensing.


\begin{center}
{\bf ADDENDUM: How to use this License for your documents\par}
\phantomsection
\end{center}

To use this License in a document you have written, include a copy of
the License in the document and put the following copyright and
license notices just after the title page:

\bigskip
\begin{quote}
	 Copyright \copyright{}  YEAR  YOUR NAME.
	 Permission is granted to copy, distribute and/or modify this document
	 under the terms of the GNU Free Documentation License, Version 1.3
	 or any later version published by the Free Software Foundation;
	 with no Invariant Sections, no Front-Cover Texts, and no Back-Cover Texts.
	 A copy of the license is included in the section entitled ``GNU
	 Free Documentation License''.
\end{quote}
\bigskip
	 
If you have Invariant Sections, Front-Cover Texts and Back-Cover Texts,
replace the ``with \dots\ Texts.'' line with this:

\bigskip
\begin{quote}
	 with the Invariant Sections being LIST THEIR TITLES, with the
	 Front-Cover Texts being LIST, and with the Back-Cover Texts being LIST.
\end{quote}
\bigskip
	 
If you have Invariant Sections without Cover Texts, or some other
combination of the three, merge those two alternatives to suit the
situation.

If your document contains nontrivial examples of program code, we
recommend releasing these examples in parallel under your choice of
free software license, such as the GNU General Public License,
to permit their use in free software.
}
				\end{multicols}
\end{document}
